\slajdid{09_3-s056}
\begin{frame}[label=09_3-s056]
\gusttekst\setlength{\abovedisplayskip}{3pt}\setlength{\belowdisplayskip}{3pt}\setlength{\abovedisplayshortskip}{2pt}\setlength{\belowdisplayshortskip}{2pt}Da degradirani naponski nivoi ne bi dalje propagirali u ostatak kola na izlaz dela koji je realizovan prolaznom logikom po pravilu se dodaje standardno CMOS logičko kolam odnosno invertor, čime se dobijaju na primer i NI, NILI logička kola itd...
\par\medskip
\begin{columns}[c,onlytextwidth]\column{.39\textwidth}\centering\begin{tikzpicture}[x=.75cm,y=.65cm,font=\sitneoznake,>=Latex]\ctikzset{bipoles/length=.6cm}\begin{scope}[shift={(0,1)}]\draw(-.35,0)--(.35,0);\draw(-.4,-.13)--(.4,-.13);\draw(-.3,0)--(-.3,.4)--(-1,.4)node[ocirc]{}node[left]{$B$};\draw(.3,0)--(.3,.4)--(1,.4);\draw(0,-.13)--(0,-.8)node[ocirc]{}node[below]{$A$};\node[above]at(0,.1){};\end{scope}\begin{scope}[shift={(0,-1)}]\draw(-.35,0)--(.35,0);\draw(-.4,-.13)--(.4,-.13);\draw(-.3,0)--(-.3,.4)--(-1,.4)node[ocirc]{}node[left]{$A$};\draw(.3,0)--(.3,.4)--(1,.4);\draw(0,-.13)--(0,-.8)node[ocirc]{}node[below]{$\overline A$};\node[above]at(0,.1){};\end{scope}
\draw(1,1.4)--(1,-.6);\draw(1,.2)--(2,.2);\draw(2,-1)--(2,1.4);
\draw(3.15,1.1)--(3.15,1.9);\draw(3,1.05)--(3,1.95);\draw(2,1.4)--(2.8,1.4)node[ocirc]{}--(3,1.4);\draw(3.15,1.85)--(3.6,1.85)--(3.6,2.4);\draw(3.15,1.15)--(3.6,1.15)--(3.6,-1);\draw(3.3,2.4)--(3.9,2.4);\node[left]at(3.2,2.4){$V_{DD}$};\node[right]at(3.8,1.4){P};
\draw(3.15,-1.4)--(3.15,-.6);\draw(3,-1.45)--(3,-.55);\draw(2,-1)--(3,-1);\draw(3.15,-.65)--(3.6,-.65);\draw(3.15,-1.35)--(3.6,-1.35)--(3.6,-2)node[ground]{};\node[right]at(3.8,-1){N};
\draw(3.6,.2)--(5,.2)node[ocirc]{}node[right]{Y};
\draw(2,2.6)--(2,3.4);\draw(2.15,2.55)--(2.15,3.45);\draw(2,3.35)--(1.55,3.35)--(1.55,3.9)node[vcc]{$V_{DD}$};\draw(2,2.65)--(1.55,2.65)--(1.55,.2);\draw(2.15,3)--(2.35,3)node[ocirc]{}--(4.5,3)--(4.5,.2);\node[left]at(1.45,3){R};
\end{tikzpicture}
\column{.59\textwidth}
Na slici je prikazana i dodatna mogućnost, odnosno dodavanje tranzistora R za restauraciju naponskih nivoa. Kada prolazna logika treba da da logičku jedinicu na izlazu invertora će biti logička nula. Tada će voditi tranzistor R i on će dopuniti parazitne kapacitivnosti do pravog napona, napona logičke jedinice odnosno $V_{DD}$. Međutim treba voditi računa o dimenzijama tranzistora R pošto prilikom prelaska izlaza sa logičke nule na logičku jedinicu treba obezbediti da tranzistori prolazne logike mogu da daju dovoljan napon logičke nule kako bi se tranzistor N zakočio. Ova situacija je praktično identična kao kod pseudo NMOS invertora koji treba da da logičku nulu na izlazu. Jednačine i odnosi su tamo definisani.
\end{columns}
\end{frame}
